\documentclass[11pt]{article}

\usepackage[T1]{fontenc}
\usepackage[utf8]{inputenc}
\usepackage{lmodern}
\usepackage[margin=1in]{geometry}
\usepackage{amsmath,amssymb,amsthm,mathtools}
\usepackage{booktabs}
\usepackage{array}
\usepackage{enumitem}
\usepackage{microtype}
\usepackage{placeins}
\usepackage{xcolor}
\usepackage{hyperref}

\hypersetup{
  colorlinks=true,
  linkcolor=blue!55!black,
  citecolor=blue!55!black,
  urlcolor=blue!55!black,
  pdftitle={Eta invariants of transverse-traceless tensor curl on three-dimensional lens spaces},
  pdfauthor={Luqman Mahmood}
}

\newtheorem{theorem}{Theorem}[section]
\newtheorem{proposition}[theorem]{Proposition}
\newtheorem{lemma}[theorem]{Lemma}
\newtheorem{corollary}[theorem]{Corollary}
\theoremstyle{definition}
\newtheorem{definition}[theorem]{Definition}
\newtheorem{remark}[theorem]{Remark}
\newtheorem{example}[theorem]{Example}

\DeclareMathOperator{\Sym}{Sym}
\DeclareMathOperator{\rank}{rank}
\newcommand{\TT}{\mathrm{TT}}
\newcommand{\Z}{\mathbb Z}
\newcommand{\R}{\mathbb R}
\newcommand{\C}{\mathbb C}
\newcommand{\ind}{\mathbf 1}
\newcommand{\e}{\mathrm e}
\newcommand{\dd}{\mathrm d}
\newcolumntype{L}[1]{>{\raggedright\arraybackslash}p{#1}}

\title{Eta invariants of transverse-traceless tensor curl\\
on three-dimensional lens spaces}
\author{Luqman Mahmood\thanks{Copyright 2026 Luqman Mahmood. This paper is licensed under Creative Commons Attribution 4.0 International (CC BY 4.0); see \url{https://creativecommons.org/licenses/by/4.0/}. Software is excluded from this grant.}}
\date{30 September 2026}

\begin{document}

\maketitle

\begin{abstract}
Let $L(p,q)$ be a round lens space with $p\geq2$ and let $A_J$ be the
tensor-curl operator obtained by compressing Homma's normalized
three-dimensional self-gradient to divergence-free trace-free symmetric
tensors of rank $J\geq1$.  We show that $A_J$ is a self-adjoint projected
elliptic pseudodifferential operator with compact resolvent, and we evaluate
its eta invariant for every unit $q$ modulo $p$.  The result is an exact
finite expression combining a classical Dedekind sum, elementary signed
residue sums, finite Fourier-aliasing corrections, and an explicit low-mode
count.  For homogeneous quotients it reduces to
\[
 \eta_J(L(p,1))
 =-\frac{J\bigl(p^2-3Jp+4J^2-2\bigr)}{3p},
 \qquad p\geq\max(2,2J-1).
\]
The threshold is sharp.  We further establish rank periodicity and the
reflection $\eta_{p-J}=-\eta_J$ when $1\pm q$ are units modulo $p$, and we
exhibit four parameterized eta-zero families at $p=2J^2-1$ (with possible
coincidences at the smallest ranks).  At $J=1$ the formula recovers the
classical odd-signature value $-4s(q,p)$ and matches the normalization in
Jeffrey's Chern--Simons spectral-flow formula; in rank two it evaluates the
spectral-asymmetry (eta) contribution to the determinant phase set aside by
Castro, Lashkari, and Maloney, not a complete one-loop phase.  The operator,
its sphere spectrum, and the equivariant character method are classical; the
contribution is the explicit signed evaluation, uniform in the rank and in
prime and composite levels.  The sphere $S^3$ (the $p=1$ case) has symmetric
TT spectrum and eta invariant zero.
\end{abstract}

\section{Introduction}

Spectral determinants on spherical quotients often retain the magnitude of a
first-order spectrum while discarding its sign.  For transverse-traceless
(TT) tensors this loses precisely the spectral asymmetry measured by an eta
invariant.  The arbitrary-rank TT representations and quotient character
machinery are classical in the heat-kernel literature
\cite{DavidGaberdielGopakumar2010}.  In rank two,
Castro--Lashkari--Maloney write down the same tensor-curl operator and its lens
space degeneracies, but explicitly set aside the sign or determinant phase
because only the norm is used \cite[eqs.~(3.3), (3.13)--(3.18) and footnote~5]
{CastroLashkariMaloney2011}.  This paper evaluates that signed difference and
its integer-rank extension.

The local starting point is Homma's three-dimensional higher-spin
self-gradient \cite{Homma2001,HommaTomihisa2021}.  It is not elliptic on the
unrestricted trace-free symmetric-tensor bundle.  Its compression to the
solenoidal range, however, has principal symbol $\pm|\xi|$.  We formulate
this compression as a projected, or Toeplitz, elliptic operator.  On a round
lens quotient its two chiral branches are
\[
 V_{n+2J}\otimes V_n,
 \qquad
 V_n\otimes V_{n+2J},
 \qquad n\geq0,
\]
with signed eigenvalues $\pm(n+J+1)$ on the unit three-sphere.  Here $V_a$
denotes the irreducible $SU(2)$ representation of highest weight $a$ and
dimension $a+1$.

Our main result is a closed formula for the eta invariant after cyclic
projection.  It applies at every level $p\geq2$ and explicitly
retains both removable denominator-singular character limits and the modular
aliasing terms that are easy to lose in a generic trigonometric calculation.
At even order the half-turn is the sole additional denominator-singular
class; its two chiral TT character products agree termwise, so its signed
contribution is zero before analytic continuation.
Three consequences are particularly compact:

\begin{enumerate}[label=(\roman*)]
\item the homogeneous polynomial in $p$ and $J$ displayed in the abstract;
\item rank periodicity and antisymmetry on nonsingular lens rows; and
\item four infinite arithmetic zero families at $p=2J^2-1$.
\end{enumerate}

At $J=1$ the result is the established curl or odd-signature eta invariant
\cite{APS1975,Bar2019}.  This fixes the normalization.  A second consistency
check comes from the appearance of the same
$J=1$ term in Jeffrey's $SU(2)$ Chern--Simons spectral-flow formula
\cite{Jeffrey1992}; Appendix~\ref{app:jeffrey} records the comparison.

The operator, sphere spectrum, quotient projection, and equivariant eta
methods are antecedent machinery.  To the best of our knowledge, the explicit
signed all-rank lens formula and its uniform treatment of composite levels
have not appeared before.  The result is a spectral statement for round lens
spaces; completing a gravitational or higher-spin determinant would require
additional sectors and choices not considered here.

Section~\ref{sec:conventions} fixes the quotient and orientation conventions.
Sections~3 and 4 construct the projected operator and continue its eta
series.  Sections~5--7 prove the closed formula and its rank-symmetry and
zero-family consequences.  The final sections record computational checks,
place the result among its antecedents, and discuss open problems.

\section{Lens-space and orientation conventions}
\label{sec:conventions}

Write the unit three-sphere as
\[
 S^3=\{(z_1,z_2)\in\C^2:|z_1|^2+|z_2|^2=1\}.
\]
For $p\geq2$ and $q\in(\Z/p\Z)^\times$, the generator of $\Z_p$ acts by
\[
 (z_1,z_2)\longmapsto
 \bigl(\e^{2\pi i/p}z_1,\e^{2\pi iq/p}z_2\bigr),
\]
and $L(p,q)$ denotes the quotient with its complex quotient orientation.  We
choose the branch $V_{n+2J}\otimes V_n$ to have positive eigenvalue.  Changing
either global convention reverses every eta value simultaneously.
Throughout, $\eta$ means the ordinary eta invariant of $A_J$, not the reduced eta
$(\eta+h)/2$; the zero-mode convention is therefore kept separate.
The round-boundary completion in \cite{Mahmood2026SteinWeissAPSBoundary}
uses this same positive-left-heavy TT convention; its auxiliary finite
character is subtracted when forming the full boundary eta invariant.
The case $p=1$ is the sphere $S^3$, whose two TT branches have equal
multiplicity at every level and hence eta invariant zero; it is not included
in the nontrivial lens-space theorem below.

Let $q^*\in\{1,\ldots,p-1\}$ be the least positive inverse of $q$ modulo
$p$.  For an integer $k$, let $\langle k\rangle_p\in\{0,\ldots,p-1\}$ be its
least nonnegative residue and define
\begin{equation}
\rho_p(k)=
\begin{cases}
0,&p\mid k,\\
p-2\langle k\rangle_p,&p\nmid k,
\end{cases}
\qquad
\delta_p(k)=\ind_{p\mid k}.
\label{eq:rho-delta}
\end{equation}
The sawtooth function and Dedekind sum are
\[
 \left(\!\left(x\right)\!\right)=
 \begin{cases}
 x-\lfloor x\rfloor-\tfrac12,&x\notin\Z,\\
 0,&x\in\Z,
 \end{cases}
 \qquad
 s(q,p)=\sum_{a=0}^{p-1}
 \left(\!\left(\frac ap\right)\!\right)
 \left(\!\left(\frac{qa}{p}\right)\!\right).
\]

For $d\geq0$, set $W_d=\{-d,-d+2,\ldots,d\}$.  For $a,b\geq1$ with
$a\equiv b\pmod2$, the cyclic invariant count in
$V_{a-1}\otimes V_{b-1}$ is
\begin{equation}
 I_{a,b}(p,q)=\#\left\{(r,t)\in W_{a-1}\times W_{b-1}:
 (1+q)r+(1-q)t\equiv0\pmod{2p}\right\}.
 \label{eq:Iab}
\end{equation}
Finally, define the finite low-mode correction
\begin{equation}
 L_J(p,q)=\sum_{m=1}^{J-1}
 \left[I_{J-m,J+m}(p,q)-I_{J+m,J-m}(p,q)\right].
 \label{eq:low-mode}
\end{equation}
The sum is empty at $J=1$.

\begin{center}
\begin{tabular}{c|l}
\hline
symbol & meaning\\
\hline
$A_J$ & normalized rank-$J$ self-gradient compressed to TT tensors\\
$\rho_p(k)$ & signed least-residue function in \eqref{eq:rho-delta}\\
$\delta_p(k)$ & divisibility indicator $\mathbf 1_{p\mid k}$\\
$I_{a,b}(p,q)$ & cyclic invariant count in
  $V_{a-1}\otimes V_{b-1}$\\
$L_J(p,q)$ & finite subtraction of formal modes below the TT threshold\\
\hline
\end{tabular}
\end{center}

\begin{remark}[Radius scaling]
\label{rem:radius-scaling}
If the round metric is rescaled from unit radius to radius $R$, then
$A_J$ and every nonzero eigenvalue scale by $R^{-1}$.  Consequently
$\eta_{A_J,R}(s)=R^s\eta_{A_J,1}(s)$, and the value
$\eta_{A_J,R}(0)$ in the theorem below is independent of $R$.
\end{remark}

The quotient classification symmetries give, in these conventions,
\begin{equation}
 \eta_J(L(p,q^*))=\eta_J(L(p,q)),
 \qquad
\eta_J(L(p,-q))=-\eta_J(L(p,q)).
\label{eq:lens-symmetries}
\end{equation}
Indeed, swapping the two complex coordinates and then replacing the cyclic
generator by its $q^*$th power gives the orientation-preserving
$q\mapsto q^*$ identification.  Conjugating the second complex coordinate
gives the orientation-reversing $q\mapsto-q$ identification; orientation
reversal negates tensor curl and hence its eta invariant.

\section{The projected TT tensor-curl operator}
\label{sec:operator}

Let $(M,g)$ be a closed oriented constant-curvature three-manifold.  We work
on the complexification of the real tensor bundle and put
\[
 E_J=\Sym^J_0T^*M.
\]
Let $D^0_{2J}$ be Homma's self-gradient on the even-spin bundle identified
with $E_J$, let
\[
 U_J=\frac{D^0_{2J}}{2J},
\]
and let $\delta_J:E_J\to E_{J-1}$ be trace-free symmetric divergence.  If
$P_J$ is the exact $L^2$-orthogonal projection onto $\ker\delta_J$, define
\begin{equation}
 A_J=P_JU_JP_J:
 P_JH^1(E_J)\subset P_JL^2(E_J)\longrightarrow P_JL^2(E_J).
 \label{eq:AJ}
\end{equation}

\begin{theorem}[Projected TT operator]
\label{thm:tt-operator}
The operator $A_J$ is self-adjoint and has compact resolvent.  Its compressed
principal symbol on $\operatorname{ran}\sigma(P_J)(\xi)$ has eigenvalues
$+|\xi|$ and $-|\xi|$ for every $\xi\neq0$.  Hence $A_J$ is elliptic in the
projected pseudodifferential calculus.  On the unit round sphere its TT
spectrum is
\begin{align}
 V_{n+2J}\otimes V_n &: +(n+J+1),\nonumber\\
 V_n\otimes V_{n+2J} &: -(n+J+1),
 \qquad n\geq0.
 \label{eq:sphere-spectrum}
\end{align}
The spectrum on a round spherical quotient is obtained by taking invariant
vectors in each representation.  In particular, there are no zero modes.
\end{theorem}

\begin{proof}
Write
\[
 d(\xi)=\sigma_1(\delta_J)(x,\xi),
 \qquad
 u(\xi)=\sigma_1(U_J)(x,\xi).
\]
The even spin module $V_{2J}$ factors through $SO(3)$ and identifies with
the fiber of $\Sym^J_0T^*M$.  After rotating a nonzero covector $\xi$ to a
fixed axis, its stabilizer decomposes that fiber into one-dimensional weight
spaces
\[
 E_{J,x}=\bigoplus_{m=-J}^{J}\C e_m .
\]
Homma's explicit symbol formulae, in the Hermitian Fourier convention and
with $U_J=D^0_{2J}/(2J)$, give
\begin{equation}
 u(\xi)e_m=\frac mJ|\xi|e_m,
 \qquad
 \|d(\xi)e_m\|^2
 =|\xi|^2\frac{J^2-m^2}{J(2J-1)}.
 \label{eq:operator-fiber-symbols}
\end{equation}
The target $E_{J-1,x}$ has precisely the weights
$-(J-1),\ldots,J-1$.  Thus $d(\xi)$ is nonzero on every source line that
maps to a target weight, vanishes exactly on the two extreme lines, and is
surjective.  Consequently
\[
 \ker d(\xi)=\operatorname{span}\{e_{-J},e_J\},
 \qquad
 u(\xi)|_{\ker d(\xi)}
 =\operatorname{diag}(-|\xi|,+|\xi|),
\]
up to the single global sign fixed in Section~\ref{sec:conventions}.

We next construct the exact global projection.  Put
\[
 \Delta_{\delta,J}=\delta_J\delta_J^*,
 \qquad
 K_{J-1}=\ker\delta_J^* .
\]
Surjectivity of $d(\xi)$ implies that $\Delta_{\delta,J}$ is elliptic of
order two.  It is nonnegative and
$\ker\Delta_{\delta,J}=K_{J-1}$.  Hence $K_{J-1}$ is finite dimensional and
smooth, and its orthogonal projector $\Pi_{K_{J-1}}$ is smoothing.  The
positive operator $\Delta_{\delta,J}+\Pi_{K_{J-1}}$ is invertible, with
\[
 G_J=(\Delta_{\delta,J}+\Pi_{K_{J-1}})^{-1}\in\Psi^{-2},
 \qquad
 \Delta_{\delta,J}G_J
 =G_J\Delta_{\delta,J}=I-\Pi_{K_{J-1}}.
\]
Since $\operatorname{ran}\delta_J\perp K_{J-1}$,
$\Pi_{K_{J-1}}\delta_J=0$.  Therefore
\[
 P_J=I-\delta_J^*G_J\delta_J
\]
satisfies
\[
 \delta_JP_J
 =\delta_J-\Delta_{\delta,J}G_J\delta_J
 =\delta_J-(I-\Pi_{K_{J-1}})\delta_J=0.
\]
It is self-adjoint and acts as the identity on $\ker\delta_J$.  It is
therefore exactly the $L^2$-orthogonal projection onto that kernel.  Its
principal symbol is
\begin{equation}
 p(\xi)=I-d(\xi)^*
 \bigl(d(\xi)d(\xi)^*\bigr)^{-1}d(\xi),
 \label{eq:solenoidal-symbol}
\end{equation}
the orthogonal projection onto $\ker d(\xi)$.  In particular
$P_J\in\Psi^0$ and is bounded on every Sobolev space; compare the standard
solenoidal construction in \cite{GouezelLefeuvre2021}.

The compressed principal symbol of $A_J=P_JU_JP_J$ is
\[
 p(\xi)u(\xi)p(\xi):
 \operatorname{ran}p(\xi)\longrightarrow\operatorname{ran}p(\xi).
\]
By \eqref{eq:operator-fiber-symbols}, its eigenvalues are
$\pm|\xi|$, so it is invertible for $\xi\neq0$.  This proves projected
ellipticity \cite{SavinSternin1999,Seiler2012}.

It remains to identify the closed domain.  Let $Q_J=I-P_J$ and
\[
 \Lambda=(1+\nabla^*\nabla)^{1/2},\qquad
 \mathcal B_J=P_JU_JP_J+Q_J\Lambda Q_J.
\]
The complex-power calculus gives $\Lambda\in\Psi^1$ with scalar principal
symbol $|\xi|$.  Because $U_J$ is formally self-adjoint and $P_J,Q_J$ are
orthogonal projections, $\mathcal B_J$ is formally self-adjoint.  Relative
to
\[
 E_{J,x}=\operatorname{ran}p(\xi)\oplus
 \operatorname{ran}(I-p(\xi)),
\]
its order-one symbol is block diagonal:
\[
 \sigma_1(\mathcal B_J)(\xi)=
 \begin{pmatrix}
 p(\xi)u(\xi)p(\xi)|_{\operatorname{ran}p(\xi)}&0\\
 0&|\xi|I
 \end{pmatrix}.
\]
Both blocks are invertible, hence $\mathcal B_J$ is an elliptic classical
operator of order one.

Let $\mathcal B_{\min}$ be its graph closure from smooth sections and
$\mathcal B_{\max}$ its distributional maximal realization.  If
$v\in L^2$ and $\mathcal B_Jv\in L^2$, an order-$-1$ parametrix $C$ with
$C\mathcal B_J=I-R$, $R$ smoothing, gives
\[
 v=C\mathcal B_Jv+Rv\in H^1.
\]
Conversely, an order-one operator maps $H^1$ continuously to $L^2$, so
$\mathcal D(\mathcal B_{\max})=H^1(E_J)$.  The elliptic estimate
\[
 \|v\|_{H^1}\leq C_0
 \bigl(\|\mathcal B_Jv\|_{L^2}+\|v\|_{L^2}\bigr)
\]
and density of smooth sections in $H^1$ show that
$\mathcal D(\mathcal B_{\min})=H^1(E_J)$ as well.  Formal symmetry now yields
\[
 \mathcal B_{\min}^*=\mathcal B_{\max}=\mathcal B_{\min};
\]
thus $\mathcal B_J$ is self-adjoint on $H^1(E_J)$.  For nonreal $z$,
elliptic regularity makes $(\mathcal B_J-z)^{-1}:L^2\to H^1$ bounded.
Composing with the compact Rellich inclusion proves that its resolvent is
compact.

Algebraically,
\[
 P_J\mathcal B_J=\mathcal B_JP_J=P_JU_JP_J.
\]
Thus $\operatorname{ran}P_J$ is reducing.  The corresponding restriction is
$A_J$, with domain
\[
 H^1(E_J)\cap\operatorname{ran}P_J=P_JH^1(E_J).
\]
It is self-adjoint, and its resolvent is the compact compression of the
resolvent of $\mathcal B_J$.

Finally, Homma's lowering operator $D^-_{2J}$ is a nonzero constant multiple
of $\delta_J$ under the even-spin/STF identification.  On constant
curvature, Homma's intertwining identity reads
\[
 D^0_{2J-2}D^-_{2J}=D^-_{2J}D^0_{2J}.
\]
Therefore $U_J$ preserves $\ker\delta_J$, and the compression agrees there
with the local self-gradient.  On $S^3$ the solenoidal trace-free sector is
\[
 \bigoplus_{n\geq0}
 \left[
 V_{n+2J}\otimes V_n
 \oplus
 V_n\otimes V_{n+2J}
 \right].
\]
Homma's signed spectrum for $D^0_{2J}$ is
$+2J(n+J+1)$ on the first summand and $-2J(n+J+1)$ on the second
\cite{Homma2001,HommaTomihisa2021}.  Dividing by $2J$ proves
\eqref{eq:sphere-spectrum}.  If a finite group acts freely by
orientation-preserving isometries, pullback identifies quotient sections
with invariant sections upstairs and intertwines both $\delta_J$ and $U_J$.
The quotient eigenspaces are therefore exactly the invariant vectors in the
displayed modules.  Their eigenvalue magnitudes are $n+J+1>0$, so no zero
mode occurs.
\end{proof}

\begin{remark}
The theorem does not make $U_J$ elliptic on all of $E_J$, and the range of
$P_J$ is not the section space of a local rank-two vector bundle.  The
transverse projection is nonlocal.  The intrinsic projected operator alone
is not asserted to be a local boundary operator on an arbitrary
three-manifold.  The separate round boundary interpretation is described in
the next subsection and is not an assumption of the eta calculation.
\end{remark}

\begin{lemma}[Rank-one normalization]
\label{lem:rank-one-normalization}
Under the metric identification $E_1=T^*M\otimes\mathbb C$, the normalized
self-gradient is
\[
 U_1=\frac{D^0_2}{2}=*d .
\]
Moreover $\delta_1$ is the codifferential up to the convention fixed above,
so $A_1$ is the classical curl operator on coclosed one-forms.  On a lens
space this is the coexact sector, and in the orientation of
Section~\ref{sec:conventions},
\[
 \eta(A_1;L(p,q))=-4s(q,p).
\]
\end{lemma}

\begin{proof}
In an oriented orthonormal frame the unique $SO(3)$-equivariant
self-gradient on one-forms has components
$(U_1\alpha)_i=\varepsilon_i{}^{jk}\nabla_j\alpha_k$, which is $*d\alpha$.
This fixes the operator normalization rather than only its spectrum.  The
lowering operator is contraction with $\nabla$, hence a
convention-fixed nonzero multiple of the codifferential; its kernel is the
coclosed subspace.  Because $H^1(L(p,q);\mathbb R)=0$, coclosed one-forms
have no harmonic summand and are coexact.  The classical
odd-signature/curl formula in the chosen quotient orientation is
$-4s(q,p)$ \cite{APS1975,Bar2019}.
\end{proof}

\subsection{Relation to the round boundary operator}
\label{sec:round-boundary-interface}

For comparison with the companion boundary problem, on a round quotient the
rank-$J$ TT sector is a reducing block on which the companion boundary operator
restricts to $A_J$ in the convention of Section~\ref{sec:conventions}.  The
square completion, twisted-Dirac identification, and assembly of the full
Atiyah--Patodi--Singer boundary correction---which involves auxiliary eta and
kernel terms beyond $\eta(A_J)$---are carried out in
\cite{Mahmood2026SteinWeissAPSBoundary} and are not needed here.  The
rankwise Dirac-to-tensor-curl transfer and the generic McKay--Chebyshev
algebra appear separately in
\cite{Mahmood2026DiracToTensorCurlEtaTransfer,
Mahmood2026BivariateMcKay}.

\section{Cyclic projection and eta continuation}
\label{sec:continuation}

From \eqref{eq:Iab} and \eqref{eq:sphere-spectrum}, put
\begin{equation}
 d^+_{J,n}=I_{n+2J+1,n+1}(p,q),\qquad
 d^-_{J,n}=I_{n+1,n+2J+1}(p,q),
 \label{eq:branch-multiplicities}
\end{equation}
and $D_J(n;p,q)=d^+_{J,n}-d^-_{J,n}$.  For $\Re s$ large, the eta function is
\begin{equation}
 \eta_{A_J}(s)=\sum_{n\geq0}
 \frac{D_J(n;p,q)}{(n+J+1)^s}.
 \label{eq:eta-series}
\end{equation}

\begin{proposition}[Affine residue continuation]
\label{prop:eta-continuation}
For each residue $0\leq r<2p$ there are integers
$A_r=A_r(J,p,q)$ and $B_r=B_r(J,p,q)$ such that
\begin{equation}
 D_J(r+2pk;p,q)=A_rk+B_r\qquad(k\geq0).
 \label{eq:affine-residue}
\end{equation}
With $a_r=(r+J+1)/(2p)$,
\begin{equation}
 \eta_{A_J}(s)=(2p)^{-s}\sum_{r=0}^{2p-1}
 \left[A_r\zeta(s-1,a_r)+(B_r-A_ra_r)\zeta(s,a_r)\right].
 \label{eq:hurwitz-continuation}
\end{equation}
In particular, the eta function is regular at $s=0$ and
\begin{equation}
 \eta_{A_J}(0)=\sum_{r=0}^{2p-1}
 \left[A_r\left(\frac{a_r^2}{2}-\frac1{12}\right)
 +B_r\left(\frac12-a_r\right)\right].
 \label{eq:eta-zero-hurwitz}
\end{equation}
\end{proposition}

\begin{proof}
For a highest weight $h\geq0$ and $a\in\Z/(2p)\Z$, let
\[
 c_h(a)=\#\{w\in W_h:w\equiv a\pmod{2p}\}.
\]
The interval of weights grows by one full residue period at each end when
$h$ is replaced by $h+2p$.  Consequently
\begin{equation}
 c_{h+2p}(a)=c_h(a)+
 2\ind_{a\equiv h\;(\mathrm{mod}\;2)}.
 \label{eq:residue-count-recurrence}
\end{equation}
Introduce the bilinear invariant-counting form
\[
 \mathcal I_{p,q}(c,d)=
 \sum_{\substack{a,b\;(\mathrm{mod}\;2p)\\
 (1+q)a+(1-q)b\equiv0\;(\mathrm{mod}\;2p)}}
 c(a)d(b).
\]
Fix $0\leq r<2p$, put $c=c_{r+2J}$ and $d=c_r$, and let
$u(a)=2\ind_{a\equiv r\;(\mathrm{mod}\;2)}$.  The two highest weights have
the same parity.  Iterating \eqref{eq:residue-count-recurrence} therefore
gives
\[
 c_{r+2J+2pk}=c+ku,\qquad c_{r+2pk}=d+ku.
\]
Hence
\begin{align*}
 D_J(r+2pk;p,q)
={}&\mathcal I_{p,q}(c+ku,d+ku)
   -\mathcal I_{p,q}(d+ku,c+ku)\\
={}&\mathcal I_{p,q}(c,d)-\mathcal I_{p,q}(d,c)\\
&+k\bigl[
 \mathcal I_{p,q}(u,d)+\mathcal I_{p,q}(c,u)
 -\mathcal I_{p,q}(u,c)-\mathcal I_{p,q}(d,u)
 \bigr].
\end{align*}
The two quadratic terms are identical and cancel.  This is
\eqref{eq:affine-residue}, with the constant and linear coefficients in the
last display equal to $B_r$ and $A_r$.

Splitting \eqref{eq:eta-series} into these residue classes, and writing
$a_r=(r+J+1)/(2p)$, gives for $\Re s$ large
\begin{align*}
 \eta_{A_J}(s)
 &=(2p)^{-s}\sum_{r=0}^{2p-1}
   \sum_{k\geq0}(A_rk+B_r)(k+a_r)^{-s}\\
 &=(2p)^{-s}\sum_{r=0}^{2p-1}
 \left[A_r\zeta(s-1,a_r)+(B_r-A_ra_r)\zeta(s,a_r)\right].
\end{align*}
The right side is the meromorphic continuation
\eqref{eq:hurwitz-continuation} and is regular at zero.  Finally,
\[
 \zeta(0,a)=\frac12-a,\qquad
 \zeta(-1,a)=-\frac12\left(a^2-a+\frac16\right),
\]
and substitution gives \eqref{eq:eta-zero-hurwitz}.
\end{proof}

The explicit Hurwitz formula proves regularity for the reducing TT spectrum
directly.  It is consistent with the general residue criteria for eta
functions of elliptic pseudodifferential operators
\cite{Ponge2003,SeilerStrohmaier2010}.

\section{The all-rank lens-space formula}
\label{sec:master}

\begin{theorem}[All-rank eta formula]
\label{thm:eta-master}
For every integer $J\geq1$, every $p\geq2$, and every
$q\in(\Z/p\Z)^\times$, the standard eta invariant of \eqref{eq:AJ} is
\begin{align}
\eta_J(L(p,q))={}&-4J\,s(q,p)\nonumber\\
&+\frac{2}{p}\sum_{a=1}^{J-1}(J-a)
  \left[\rho_p(aq)+\rho_p(aq^*)\right]\nonumber\\
&-2\!\!\sum_{\substack{a,b\geq1\\a+b\leq J-1}}
  (J-a-b)\left[\delta_p(aq-b)-\delta_p(aq+b)\right]\nonumber\\
&-L_J(p,q),
\label{eq:master-formula}
\end{align}
where \eqref{eq:rho-delta} and \eqref{eq:low-mode} define every term.
\end{theorem}

The formula is valid without a parity or primality assumption.  The third
line records finite Fourier-aliasing congruences, including composite-level
coincidences, and the fourth line removes the formal modes below the physical
threshold $n+J+1=J+1$.

\subsection{Character reduction}

For a nonidentity group element indexed by $1\leq\ell<p$, set
\[
 \alpha_\ell=\frac{\pi\ell q}{p},\qquad
 \beta_\ell=\frac{\pi\ell}{p}.
\]
Also put
\[
 u_\ell=\alpha_\ell+\beta_\ell,\qquad
 v_\ell=\beta_\ell-\alpha_\ell,
\]
and write
\[
 \chi_h(\theta)=\frac{\sin((h+1)\theta)}{\sin\theta}
 =\sum_{w\in W_h}e^{iw\theta}.
\]
Extend the Weyl formula to integer indices by
\[
 \chi_{-1}=0,\qquad \chi_{-h-2}=-\chi_h\quad(h\geq0).
\]
\begin{lemma}[Central descent]
\label{lem:central-descent}
The cyclic action on each TT branch and on every invariant count
$I_{a,b}(p,q)$ in \eqref{eq:Iab} is independent of the lift from $SO(4)$ to
$Spin(4)=SU(2)_L\times SU(2)_R$, for both odd and even $p$.
\end{lemma}

\begin{proof}
The two lifts differ by the diagonal central element $(-1,-1)$.  On
$V_{n+2J}\otimes V_n$ it acts by
\[
 (-1)^{(n+2J)+n}=(-1)^{2(n+J)}=1,
\]
and the exchanged TT branch has the same exponent.  Likewise,
$V_{a-1}\otimes V_{b-1}$ in \eqref{eq:Iab} has central exponent $a+b-2$,
which is even because $a\equiv b\pmod2$.  Thus all representations entering
the signed projector descend to $SO(4)$ and the order-$p$ cyclic average is
well-defined without a parity assumption on $p$.
\end{proof}

For $a\equiv b\pmod2$, character averaging gives
\begin{equation}
 I_{a,b}(p,q)=\frac1p\sum_{\ell=0}^{p-1}
 \chi_{a-1}(u_\ell)\chi_{b-1}(v_\ell).
 \label{eq:character-projector}
\end{equation}
Indeed, expansion produces the phase
$\exp\{\pi i\ell[(1+q)r+(1-q)t]/p\}$, and its cyclic average is exactly the
congruence in \eqref{eq:Iab}.

For $m\geq1$, define
\begin{align*}
 D_J^{\mathrm{form}}(m;\ell)
={}&\chi_{m+J-1}(u_\ell)\chi_{m-J-1}(v_\ell)\\
&-\chi_{m-J-1}(u_\ell)\chi_{m+J-1}(v_\ell).
\end{align*}

\begin{lemma}[The even half-turn]
\label{lem:even-half-turn}
If $p$ is even, then the unique half-turn $\ell=p/2$ has
\[
 D_J^{\mathrm{form}}(m;p/2)=0
 \qquad(m\in\Z).
\]
Thus it contributes no signed eta term, before analytic continuation.
\end{lemma}

\begin{proof}
Because $q$ is a unit modulo even $p$, it is odd.  Hence at $\ell=p/2$,
\[
 u_\ell=\frac{\pi(1+q)}2,
 \qquad
 v_\ell=\frac{\pi(1-q)}2
\]
are both integral multiples of $\pi$.  For $\varepsilon=\pm1$ the central
character value is $\chi_h(\varepsilon)=(h+1)\varepsilon^h$.  Put
$A=m+J-1$, $B=m-J-1$, and let $\varepsilon_u,\varepsilon_v$ be the two
central signs.  Since $A-B=2J$,
\[
 \frac{\varepsilon_u^A\varepsilon_v^B}
      {\varepsilon_u^B\varepsilon_v^A}
 =\left(\frac{\varepsilon_u}{\varepsilon_v}\right)^{2J}=1.
\]
The common amplitude is $(A+1)(B+1)$, so exchanging the two indices leaves
the character product unchanged, including the extended-character cases.
Their signed difference is zero.
\end{proof}

The identity element contributes zero.  Product-to-sum gives
\begin{equation}
 D_J^{\mathrm{form}}(m;\ell)=
 \frac{
 \sin(2m\beta_\ell)\sin(2J\alpha_\ell)
 -\sin(2m\alpha_\ell)\sin(2J\beta_\ell)}
 {\sin(\alpha_\ell+\beta_\ell)\sin(\beta_\ell-\alpha_\ell)}.
\label{eq:formal-character}
\end{equation}
At a vanishing denominator this expression means the removable character
limit in the preceding definition.

For $\Re s$ sufficiently large, define the formal Dirichlet series
\[
 \eta_J^{\mathrm{form}}(s)
 =\frac1p\sum_{\ell=1}^{p-1}\sum_{m\geq1}
  D_J^{\mathrm{form}}(m;\ell)m^{-s}.
\]

\begin{lemma}[Abel and eta continuation]
\label{lem:abel-eta}
Write $\eta_J^{\mathrm{form}}:=\eta_J^{\mathrm{form}}(0)$ for the value at
zero of the analytically continued formal eta series.  This value equals its
Abel value:
\begin{equation}
 \eta_J^{\mathrm{form}}
 =\frac1{2p}\sum_{\ell=1}^{p-1}
 \frac{
 \sin(2J\alpha_\ell)\cot\beta_\ell
 -\sin(2J\beta_\ell)\cot\alpha_\ell}
 {\sin(\alpha_\ell+\beta_\ell)
  \sin(\beta_\ell-\alpha_\ell)}.
 \label{eq:formal-abel-eta}
\end{equation}
For even $p$, the summand at $\ell=p/2$ is understood through the character
definition and is zero by Lemma~\ref{lem:even-half-turn}.
\end{lemma}

\begin{proof}
For $\theta\notin\pi\Z$,
\[
 \lim_{t\uparrow1}\sum_{m\geq1}t^m\sin(2m\theta)
 =\lim_{t\uparrow1}
 \frac{t\sin(2\theta)}{1-2t\cos(2\theta)+t^2}
 =\frac12\cot\theta.
\]
This proves the formula when both character denominators are nonzero.  For
fixed $\ell$, a removable character limit is a finite linear combination of
$m^\nu\omega^m$, with $\nu\leq1$ and $\omega$ a nontrivial $2p$th root of
unity.  Here is the denominator classification explicitly.  One has
\[
 \sin u_\ell=0\Longleftrightarrow
 p\mid(q+1)\ell,\qquad
 \sin v_\ell=0\Longleftrightarrow
 p\mid(q-1)\ell.
\]
If both vanish, adding and subtracting the congruences gives
$p\mid2\ell$.  Since $1\leq\ell<p$, this is possible exactly when $p$ is
even and $\ell=p/2$; the resulting signed character is identically zero by
Lemma~\ref{lem:even-half-turn}.  If exactly one denominator vanishes, taking
the character limit differentiates only a finite exponential sum and hence
adds at most one factor of $m$.  Its $m$-dependent roots are among
$e^{\pm2i\alpha_\ell}$ and $e^{\pm2i\beta_\ell}$.  The latter are nontrivial
because $0<\ell<p$, and the former are nontrivial because $q$ is a unit
modulo $p$.  Thus $\omega\neq1$ in every remaining removable case, and
\[
 \sum_{m\geq1}^{\mathrm{Abel}}\omega^m=\operatorname{Li}_0(\omega),
 \qquad
 \sum_{m\geq1}^{\mathrm{Abel}}m\omega^m
 =\operatorname{Li}_{-1}(\omega).
\]
These are also the values at $s=0$ of the corresponding continued Dirichlet
series, so Abel and eta continuation agree in every removable case.
\end{proof}

If
\[
 x=\cos(2\alpha_\ell),\quad y=\cos(2\beta_\ell),\quad
 f_J(x)=(1-x)U_{J-1}(x),
\]
then
\[
 \sin(2J\alpha_\ell)=\cot\alpha_\ell\,f_J(x),\qquad
 \sin(\alpha_\ell+\beta_\ell)
 \sin(\beta_\ell-\alpha_\ell)=\frac{x-y}{2}.
\]
Consequently
\begin{equation}
 \eta_J^{\mathrm{form}}
 =\frac1p\sum_{\ell=1}^{p-1}
 \cot\alpha_\ell\cot\beta_\ell\,P_J(x,y),
 \qquad
 P_J(x,y)=\frac{f_J(x)-f_J(y)}{x-y}.
 \label{eq:formal-polynomial-kernel}
\end{equation}
At $x=y$ we set $P_J(x,x)=f_J'(x)$.  This polynomial continuation includes
$q=\pm1$ and composite singular group elements without division by zero.

\begin{lemma}[All-rank generating function]
\label{lem:chebyshev}
The divided differences satisfy
\begin{equation}
 \sum_{J\geq1}P_J(x,y)z^J
 =-\frac{z(1-z)^2}
 {(1-2xz+z^2)(1-2yz+z^2)}.
 \label{eq:cheb-generating}
\end{equation}
Equivalently,
\[
 \frac{U_n(x)-U_n(y)}{x-y}
 =2\sum_{a=0}^{n-1}U_a(x)U_{n-1-a}(y).
\]
\end{lemma}

\begin{proof}
Insert the standard generating function
$\sum_{n\geq0}U_n(x)z^n=(1-2xz+z^2)^{-1}$ and simplify the divided
difference.  The coefficient identity follows by comparing powers of $z$.
\end{proof}

\begin{lemma}[Exact trigonometric decomposition]
\label{lem:compact-trig}
For angles at which the cotangents are defined,
\begin{align}
&\cot A\cot B\,P_J(\cos2A,\cos2B)\nonumber\\
={}&-J\cot A\cot B\nonumber\\
&+2\sum_{r=1}^{J-1}(J-r)
\left[\sin(2rA)\cot B+\sin(2rB)\cot A\right]\nonumber\\
&-4\!\!\sum_{\substack{r,s\geq1\\r+s\leq J-1}}
(J-r-s)\sin(2rA)\sin(2sB).
\label{eq:compact-trig-decomposition}
\end{align}
\end{lemma}

\begin{proof}
Assume first that both cotangents are nonzero.  Let
\[
 D_A=1-2z\cos2A+z^2,\qquad
 D_B=1-2z\cos2B+z^2.
\]
The generating function of the right side is
\begin{align*}
 &-\frac{z}{(1-z)^2}\cot A\cot B\\
 &+\frac{2z^2}{(1-z)^2}
 \left[\frac{\sin2A\cot B}{D_A}
      +\frac{\sin2B\cot A}{D_B}\right]
 -\frac{4z^3\sin2A\sin2B}{(1-z)^2D_AD_B}.
\end{align*}
Put the terms over a common denominator and use
$\sin2A/\cot A=2\sin^2A$ together with
$D_A=(1-z)^2+4z\sin^2A$, and similarly for $B$.  The numerator becomes
$-z(1-z)^4\cot A\cot B$.  After cancelling $(1-z)^2$, the result is
\[
 -\frac{z(1-z)^2\cot A\cot B}{D_AD_B},
\]
which is $\cot A\cot B$ times
\eqref{eq:cheb-generating}.  Equality of coefficients proves the lemma.
Both sides of \eqref{eq:compact-trig-decomposition} are continuous at angles
where a cotangent vanishes, so the identity extends to those cases by
continuity.
\end{proof}

\begin{lemma}[Finite Fourier identities]
\label{lem:finite-fourier}
For all integers $k,u,v$,
\begin{align}
 \frac1p\sum_{\ell=1}^{p-1}
 \cot\alpha_\ell\cot\beta_\ell&=4s(q,p),
 \label{eq:cot-dedekind}\\
 \sum_{\ell=1}^{p-1}
 \sin\left(\frac{2\pi k\ell}{p}\right)
 \cot\left(\frac{\pi\ell}{p}\right)&=\rho_p(k),
 \label{eq:sine-cot-rho}\\
 \frac1p\sum_{\ell=1}^{p-1}
 \sin\left(\frac{2\pi u\ell}{p}\right)
 \sin\left(\frac{2\pi v\ell}{p}\right)
 &=\frac12\left[\delta_p(u-v)-\delta_p(u+v)\right].
 \label{eq:sine-orthogonality}
\end{align}
The missing factor $1/p$ in \eqref{eq:sine-cot-rho} is intentional:
$\rho_p(k)$ is the unnormalized finite sine--cotangent sum.
\end{lemma}

\begin{proof}
Equation \eqref{eq:cot-dedekind} is the classical cotangent expression for
the Dedekind sum, after permuting nonzero residues by $q$.  If
$1\leq k<p$, then
\[
 \sin(2kx)\cot x
 =1+\cos(2kx)+2\sum_{a=1}^{k-1}\cos(2ax).
\]
Summing at $x=\pi\ell/p$ gives $p-2k$; reducing an arbitrary $k$ modulo
$p$ proves \eqref{eq:sine-cot-rho}.  Product-to-sum and finite exponential
orthogonality give \eqref{eq:sine-orthogonality}.
\end{proof}

\begin{proof}[Proof of Theorem~\ref{thm:eta-master}]
Apply Lemma~\ref{lem:compact-trig} to
\eqref{eq:formal-polynomial-kernel} and average over $\ell$.  The first term
is $-4J s(q,p)$ by \eqref{eq:cot-dedekind}.  The first sine--cotangent term
gives $\rho_p(rq)$ by \eqref{eq:sine-cot-rho}.  In the second, the change of
variables $h\equiv q\ell\pmod p$ gives
\[
 \sum_{\ell=1}^{p-1}\sin(2r\beta_\ell)\cot\alpha_\ell
 =\rho_p(rq^*).
\]
The sine--sine terms give, by \eqref{eq:sine-orthogonality},
\[
 -2\!\!\sum_{\substack{r,s\geq1\\r+s\leq J-1}}
 (J-r-s)
 \left[\delta_p(rq-s)-\delta_p(rq+s)\right].
\]
Thus the first three lines of \eqref{eq:master-formula} equal the formal eta
value.  The modular indicators arise from finite sine orthogonality, not
only from degenerate character denominators.

These finite Fourier identities hold for every modulus $p\geq2$ and use no
inverse of $2$ modulo $p$.  When $p$ is even, the half-turn has
$\alpha_{p/2}=\pi q/2$ and $\beta_{p/2}=\pi/2$.  Hence its cotangents and all
the sine factors appearing in Lemma~\ref{lem:compact-trig} vanish; this is
the zero character contribution already identified in
Lemma~\ref{lem:even-half-turn}.  Reinstating that zero term therefore gives
the same Dedekind, residue, and modular-indicator expressions with no extra
even-level correction.

For $m\geq J+1$, \eqref{eq:character-projector} identifies the averaged
formal character with the physical signed multiplicity at $n=m-J-1$.  At
$m=J$ it vanishes because $\chi_{-1}=0$.  If $1\leq m<J$, then
\[
 \chi_{m-J-1}=-\chi_{J-m-1},
\]
so the averaged formal multiplicity is
\[
 I_{J-m,J+m}(p,q)-I_{J+m,J-m}(p,q).
\]
The physical spectrum begins at $m=J+1$.  Its eta value is therefore the
formal value minus the sum of these sub-threshold terms, namely
$L_J(p,q)$.  This gives the final line of
\eqref{eq:master-formula} with the asserted sign.
\end{proof}

The logical distinction between the formal character device and the
geometric spectrum can be summarized as
\begin{equation}
\boxed{
\begin{gathered}
\text{extended-character series}
\ \longrightarrow\
\text{Abel/Hurwitz continuation}
\ \longrightarrow\
\text{finite Fourier evaluation}\\
\longrightarrow\
\text{subtract }L_J
\ =\
\text{physical TT eta}
\end{gathered}}
\label{eq:formal-to-physical}
\end{equation}
The first three objects include formal sub-threshold indices; only the final
subtraction imposes the lower boundary $n=0$ of
\eqref{eq:sphere-spectrum}.

\begin{corollary}[Rank-one control]
\label{cor:rank-one}
At $J=1$ all correction sums vanish and
\begin{equation}
 \eta_1(L(p,q))=-4s(q,p).
 \label{eq:rank-one}
\end{equation}
This is the classical odd-signature or curl eta value in the orientation of
Section~\ref{sec:conventions}.
\end{corollary}

\begin{corollary}[Even-order anchor]
\label{cor:rp3}
For every integer $J\geq1$,
\[
 \eta_J(L(2,1))=\eta_J(\mathbb{RP}^3)=0.
\]
\end{corollary}

\begin{proof}
Modulo $2$ one has $-1\equiv1$.  The orientation identity
\eqref{eq:lens-symmetries} therefore gives
$\eta_J(L(2,1))=-\eta_J(L(2,1))$.
\end{proof}

\begin{corollary}[Rank-two specialization]
\label{cor:rank-two}
At $J=2$,
\begin{equation}
 \eta_2(L(p,q))=-8s(q,p)
 +\frac2p\bigl[\rho_p(q)+\rho_p(q^*)\bigr]-L_2(p,q).
 \label{eq:rank-two}
\end{equation}
In particular,
\[
 \eta_2(L(p,1))=-\frac{2(p^2-6p+14)}{3p}
 \qquad(p\geq3).
\]
The remaining case $p=2$ is the zero value of
Corollary~\ref{cor:rp3}.
\end{corollary}

\begin{example}[$J=2$ on $L(5,2)$]
\label{ex:J2-L52}
Here $q^*=3\equiv-2\pmod5$.  Inversion and oddness of the Dedekind sum give
$s(2,5)=s(3,5)=-s(2,5)=0$.  Moreover
\[
 \rho_5(2)=1,\qquad \rho_5(3)=-1.
\]
The sole low-mode term is
\[
 L_2(5,2)=I_{1,3}(5,2)-I_{3,1}(5,2)=1-1=0:
\]
in each count only the zero weight satisfies the congruence.  Thus every
piece of \eqref{eq:rank-two} is visible and
\[
 \eta_2(L(5,2))=0.
\]
This example also illustrates why the low-mode term must be evaluated even
when it happens to cancel.
\end{example}

\section{Homogeneous and rank-symmetry corollaries}
\label{sec:corollaries}

\begin{theorem}[Homogeneous lens spaces]
\label{cor:eta-homogeneous}
For every integer $J\geq1$ and
$p\geq\max(2,2J-1)$,
\begin{equation}
 \eta_J(L(p,1))
 =-\frac{J\left(p^2-3Jp+4J^2-2\right)}{3p}.
 \label{eq:homogeneous}
\end{equation}
For $J\geq2$ the threshold is sharp.  At the first even level below it,
$p=2J-2$, the exact value exceeds the stable expression by $2$; for $J\geq3$,
at the first odd level below it, $p=2J-3$, the excess is $4$.
\end{theorem}

\begin{proof}
Specialize \eqref{eq:master-formula} at $q=q^*=1$.  Direct summation of the
sawtooth definition gives
\[
 s(1,p)=\frac{(p-1)(p-2)}{12p},
\]
so the Dedekind term is $-J(p^2-3p+2)/(3p)$.  Since
$1\leq a\leq J-1<p$, one has $\rho_p(a)=p-2a$, and therefore
\begin{equation}
 \frac2p\sum_{a=1}^{J-1}(J-a)\,2\rho_p(a)
 =2J(J-1)-\frac{4J(J^2-1)}{3p}.
 \label{eq:homogeneous-rho}
\end{equation}

In the triangular correction, $|a-b|<p$ and $0<a+b<p$.  Thus
$\delta_p(a-b)=\ind_{a=b}$ and $\delta_p(a+b)=0$, and its value is
\begin{equation}
 -2\sum_{a=1}^{\lfloor(J-1)/2\rfloor}(J-2a)
 =\begin{cases}
   -J(J-2)/2,&J\text{ even},\\
   -(J-1)^2/2,&J\text{ odd}.
  \end{cases}
 \label{eq:homogeneous-delta}
\end{equation}

It remains to evaluate the low-mode term rather than assume it is stable.
For $q=1$, the congruence in \eqref{eq:Iab} is $r\equiv0\pmod p$.
Every weight occurring there has absolute value at most $2J-2<p$, so only
$r=0$ is possible.  It occurs precisely when the first index is odd, whence
\[
 I_{a,b}(p,1)=b\ind_{a\text{ odd}},
 \qquad
 L_J(p,1)=\sum_{\substack{1\leq m<J\\J-m\text{ odd}}}2m
 =\left\lfloor\frac{J^2}{2}\right\rfloor.
\]
Combining this with \eqref{eq:homogeneous-delta}, the last two lines of the
master formula contribute $-J(J-1)$ for either parity of $J$.  Adding this
to the Dedekind term and \eqref{eq:homogeneous-rho} yields
\[
 \eta_J(L(p,1))
 =-\frac{J(p^2-3p+2)}{3p}+J(J-1)
   -\frac{4J(J^2-1)}{3p},
\]
which is \eqref{eq:homogeneous}.

For the first even level below the stable range, let $J\geq2$ and
$p=2J-2$.  The only new low-mode event occurs at $m=J-1$.  The first factor
in $I_{2J-1,1}(p,1)$ now admits the three weights $r=-p,0,p$ rather than only
$r=0$, while $I_{1,2J-1}(p,1)$ is unchanged.  Consequently the exact
$L_J$ is smaller than its stable value by $2$, and the exact eta invariant
exceeds \eqref{eq:homogeneous} by $2$.

Finally let $J\geq3$ and $p=2J-3$.  The preceding Dedekind, residue, and
triangular calculations remain valid.  In the low-mode count, however, the
weights $r=\pm p$ first enter, and only for $m=J-2$.  More explicitly,
\[
 I_{2,2J-2}(p,1)=0,
 \qquad
 I_{2J-2,2}(p,1)=4.
\]
Thus the exact $L_J$ is smaller than its stable value by $4$, so the exact
eta invariant exceeds the polynomial in \eqref{eq:homogeneous} by $4$.
This proves the asserted sharpness.  Notice that no primality assumption was
used.
\end{proof}

As a fixed worked example, we use the same level $137$ as the companion
coexact-mode calculation \cite{Mahmood2026CoexactHydrodynamics}.  The first
five ranks at $L(137,1)$ are
\begin{equation}
 \bigl(\eta_1,\ldots,\eta_5\bigr)
 =-\frac1{137}(6120,11974,17570,22916,28020).
 \label{eq:level137-table}
\end{equation}
Equation~\eqref{eq:homogeneous} holds at every level in its stated stable
range.

\begin{theorem}[Nonsingular rank symmetries]
\label{cor:rank-symmetry}
Suppose
\[
 \gcd(1+q,p)=\gcd(1-q,p)=1.
\]
Then
\begin{equation}
 \eta_{J+p}(L(p,q))=\eta_J(L(p,q)).
 \label{eq:rank-periodicity}
\end{equation}
For $1\leq J\leq p-1$,
\begin{equation}
 \eta_{p-J}(L(p,q))=-\eta_J(L(p,q)),
 \qquad
 \eta_p(L(p,q))=0.
 \label{eq:rank-reflection}
\end{equation}
\end{theorem}

\begin{proof}
Under the stated hypothesis the denominators in
\eqref{eq:formal-character} are nonzero.  Let $F_J$ denote the formal Abel
value in \eqref{eq:formal-abel-eta}, and set
\begin{equation}
 K_J(m)=\frac1p\sum_{\ell=1}^{p-1}
 \frac{
  \sin(2m\beta_\ell)\sin(2J\alpha_\ell)
  -\sin(2m\alpha_\ell)\sin(2J\beta_\ell)}
 {\sin(\alpha_\ell+\beta_\ell)
  \sin(\beta_\ell-\alpha_\ell)}.
 \label{eq:rank-correction-kernel}
\end{equation}
The proof of Theorem~\ref{thm:eta-master} shows that
\begin{equation}
 \eta_J=F_J-c_J,
 \qquad c_J=\sum_{m=1}^{J}K_J(m),
 \label{eq:rank-formal-minus-correction}
\end{equation}
where including $m=J$ is harmless because $K_J(J)=0$.

The elementary sine identities in the rank variable give
\[
 F_{J+p}=F_J,\qquad F_{p-J}=-F_J,\qquad F_p=0.
\]
The same identities applied to \eqref{eq:rank-correction-kernel} give
\begin{align*}
 K_{J+p}(m)&=K_J(m),&K_{p-J}(m)&=-K_J(m),&K_p(m)&=0,\\
 K_J(m+p)&=K_J(m),&K_J(p-m)&=-K_J(m),&
 K_J(0)&=K_J(J)=0.
\end{align*}
Consequently a complete period in $m$ sums to zero, and hence
\[
 c_{J+p}=c_J,\qquad c_p=0.
\]
For reflection, use the vanishing of a full period and
$K_J(p-m)=-K_J(m)$:
\[
 \sum_{m=1}^{p-J}K_J(m)
 =-\sum_{m=p-J+1}^{p}K_J(m)
 =\sum_{n=0}^{J-1}K_J(n)
 =\sum_{n=1}^{J-1}K_J(n)=c_J.
\]
Therefore $c_{p-J}=-c_J$ for every $1\leq J\leq p-1$, including the
upper half of the fundamental rank range.
Substitution in \eqref{eq:rank-formal-minus-correction} proves
\eqref{eq:rank-periodicity} and \eqref{eq:rank-reflection}.
\end{proof}

\begin{remark}
The nonsingularity hypothesis has nonempty scope exactly when
$\gcd(p,6)=1$.  If $2\mid p$, every unit $q$ is odd, so both $1+q$ and
$1-q$ share the factor $2$ with $p$.  If $3\mid p$, every unit is congruent
to $\pm1\pmod 3$, so one of $1+q$ and $1-q$ shares the factor $3$ with $p$.
Conversely, if $\gcd(p,6)=1$, then $q=2$ is a unit modulo $p$ and satisfies
both gcd conditions.  Thus this theorem has nonempty scope precisely at
levels coprime to $6$, even though Theorem~\ref{thm:eta-master} holds for all
$p\geq2$.  These are arithmetic symmetries of the rank-indexed eta sequence,
not a claim of periodic physical spin.
\end{remark}

\section{Four infinite eta-zero families}
\label{sec:zeros}

\begin{theorem}[Arithmetic zero families]
\label{thm:zero-families}
Let $J\geq2$ and $p_J=2J^2-1$.  Then
\begin{align}
 \eta_J(L(p_J,J))&=0, &
 \eta_J(L(p_J,J+1))&=0,\nonumber\\
 \eta_J(L(p_J,2J-1))&=0, &
 \eta_J(L(p_J,J(J-1)))&=0.
 \label{eq:four-zero-families}
\end{align}
Their images under $q\mapsto q^{-1}$ and $q\mapsto-q$ are also zeros.
Every level $p_J$ in these families is odd.
The four displayed formulae need not represent four distinct lens orbits at
small $J$.
\end{theorem}

\begin{proof}
Write $p=p_J$.  The identities
\begin{align*}
 J(2J)&=p+1,\\
 (J+1)(p-2J+2)&\equiv1\pmod p,\\
 (2J-1)(2J+1)&=2p+1,\\
 (4J+2)J(J-1)&=(2J-1)p-1
\end{align*}
show that all four displayed residues are units and give the inverses used
below.  Define the residue contribution
\[
 R_J(q)=\frac2p\sum_{a=1}^{J-1}(J-a)
  [\rho_p(aq)+\rho_p(aq^*)].
\]

For $q=J,J+1,2J-1$, repeated Dedekind reciprocity gives, respectively,
\begin{align}
 s(J,p)&=\frac{(J-1)(J^2-J-1)}{2p},
 \label{eq:zero-dedekind-one}\\
 s(J+1,p)&=\frac{(J-3)(J-1)(J+1)}{6p},
 \label{eq:zero-dedekind-two}\\
 s(2J-1,p)&=\frac{(J-1)(2J^2-4J-3)}{6p}.
 \label{eq:zero-dedekind-three}
\end{align}
Their inverse residues are $2J$, $p-2J+2$, and $2J+1$.  No residue wraps
except the explicitly reduced second member of the middle pair, so for
$1\leq a<J$,
\begin{align*}
 \rho_p(aJ)+\rho_p(2aJ)&=2p-6Ja,\\
 \rho_p(a(J+1))+\rho_p(a(p-2J+2))&=2a(J-3),\\
 \rho_p(a(2J-1))+\rho_p(a(2J+1))&=2p-8Ja.
\end{align*}
Using
$\sum_{a=1}^{J-1}(J-a)=J(J-1)/2$ and
$\sum_{a=1}^{J-1}a(J-a)=J(J-1)(J+1)/6$ now gives
\begin{equation}
 R_J(q)=4J\,s(q,p)
 \qquad(q=J,J+1,2J-1).
 \label{eq:first-three-residue-cancellation}
\end{equation}

The remaining two master-formula terms vanish for these three choices.
Indeed, whenever $a,b\geq1$ and $a+b\leq J-1$, one has
\[
 0<aq-b<aq+b<p,
\]
so the triangular $\delta_p$ sum is zero.  For the low-mode term put
$u=(r-t)/2$ and $v=(r+t)/2$.  The two weights have the same parity, so
$u,v\in\Z$, $|u|,|v|\leq J-1$, and the congruence in \eqref{eq:Iab} becomes
\[
 v+qu\equiv0\pmod p.
\]
For each of the three $q$ values, $|v+qu|<p$ and $q>J-1$; hence
$u=v=0$.  This zero-weight solution occurs symmetrically in the exchanged
counts, so $L_J(p,q)=0$.  Equations
\eqref{eq:zero-dedekind-one}--\eqref{eq:first-three-residue-cancellation}
therefore prove the first three zeros.

For the fourth family use $q_0=4J+2$.  Its least positive inverse is
\[
 q_0^*=p-J(J-1)=J^2+J-1,
\]
and $J(J-1)\equiv-q_0^*\pmod p$.  Set
$h=\lfloor(J-1)/2\rfloor$.  For $1\leq a<J$, direct reduction gives
\begin{align*}
 \langle aq_0\rangle_p
 &=\begin{cases}aq_0,&a\leq h,\\aq_0-p,&a>h,
   \end{cases}\\
 \langle aq_0^*\rangle_p
 &=\begin{cases}
   a(2J-1)/2,&a\text{ even},\\
   [p+a(2J-1)]/2,&a\text{ odd}.
   \end{cases}
\end{align*}
Substitution into the finite weighted sum defining $R_J$ and evaluation of
the resulting arithmetic sums yields the following expression.  Explicitly,
\[
 \rho_p(aq_0)+\rho_p(aq_0^*)
 =2p-a(10J+3)+2p\ind_{a>h}-p\ind_{a\text{ odd}},
\]
and the evaluation uses
\begin{align*}
 \sum_{a=1}^{J-1}(J-a)&=\frac{J(J-1)}2,&
 \sum_{a=1}^{J-1}a(J-a)&=\frac{J(J-1)(J+1)}6,\\
 \sum_{h<a<J}(J-a)&=\frac{H(H+1)}2,\qquad H=\left\lfloor\frac J2\right\rfloor,&
 \sum_{\substack{1\leq a<J\\a\text{ odd}}}(J-a)
 &=\begin{cases}H^2,&J=2H,\\H(H+1),&J=2H+1.\end{cases}
\end{align*}
Consequently
\begin{equation}
 R_J(q_0)=
 \begin{cases}
 \displaystyle\frac{J(2J^3-9J^2+4J+6)}{3p},&J\text{ even},\\[5pt]
 \displaystyle\frac{J(J-1)(2J^2-13J-9)}{3p},&J\text{ odd}.
 \end{cases}
 \label{eq:fourth-residue-sum}
\end{equation}
The Euclidean remainder chains
\[
 \begin{array}{ll}
 p\equiv J\pmod{q_0},\quad q_0\equiv2\pmod J,
   &J\text{ odd},\\
 p\equiv-(J+1)\pmod{q_0},\quad q_0\equiv-2\pmod{J+1},
   &J\text{ even},
 \end{array}
\]
together with Dedekind reciprocity and $s(-a,b)=-s(a,b)$ give
\begin{equation}
 s(q_0,p)=
 \begin{cases}
 \displaystyle\frac{2J^3-9J^2+4J+6}{12p},&J\text{ even},\\[5pt]
 \displaystyle\frac{(J-1)(2J^2-13J-9)}{12p},&J\text{ odd}.
 \end{cases}
 \label{eq:fourth-dedekind-sum}
\end{equation}
Thus \eqref{eq:fourth-residue-sum} is exactly $4J$ times
\eqref{eq:fourth-dedekind-sum}.

It remains only to exclude corrections.  The least absolute residue of $p$
modulo $q_0$ is $J$ for odd $J$ and $J+1$ for even $J$.  In the triangular
range, $0<aq_0\pm b<2p$ and $|b|\leq J-2$; divisibility by $p$ would force
$b\equiv\pm p\pmod{q_0}$, contradicting that least-residue bound.  Hence the
$\delta_p$ term is zero.  In the low-mode range the same variables $u,v$
satisfy $|u|,|v|\leq J-1$ and $|v+q_0u|<2p$.  The values $\pm p$ are again
excluded by the least-residue bound, while $v+q_0u=0$ forces $u=v=0$.
The surviving zero weight is symmetric, so $L_J(p,q_0)=0$.  Therefore
$\eta_J(L(p,q_0))=0$, and \eqref{eq:lens-symmetries} transfers this zero to
$q=J(J-1)$ and to every stated orbit image.
\end{proof}

A complete classification of zero orbits beyond the four constructions
above and the orientation-forced condition $q^2\equiv-1\pmod p$ remains
an open arithmetic problem.

\section{Computational verification}
\label{sec:computational-verification}

Exact computations were used to test the projector, Hurwitz-zeta,
Dedekind--Rademacher, and compact-residue forms of the answer across prime,
composite, denominator-singular, and sharp-boundary cases.  Exact symbolic
fiber matrices were compared through rank $J=10$, and independent
weight-space calculations tested the resulting formulas through rank $J=64$.
Separate checks covered the even half-turn, the $\mathbb{RP}^3$ anchor, and
the four zero families.  A companion data set contains $16{,}177$ exact
rank-one spectral-flow rows and higher-rank regression values.  All
comparisons use rational or integer arithmetic; they test the formulas and
conventions but are not substitutes for the proofs.

\FloatBarrier

\section{Relation to antecedents and scope}
\label{sec:scope}

Homma supplies the higher-spin self-gradient, its signed sphere branches, and
the constant-curvature TT preservation identity
\cite{Homma2001,HommaTomihisa2021}.  David--Gaberdiel--Gopakumar supply the
arbitrary-rank TT representations and quotient character machinery
\cite{DavidGaberdielGopakumar2010}.  B\"ar and APS provide the rank-one curl or
odd-signature control \cite{Bar2019,APS1975}.  In rank two,
Castro--Lashkari--Maloney give the direct operator and explicit unresolved
phase gap \cite{CastroLashkariMaloney2011}.  Basu gives unsigned spin-three
and spin-four lens-space determinant expressions and conjectures the pattern
above spin four \cite{Basu2015}.  Degeratu's Molien formula evaluates twisted Dirac eta
invariants, not the projected TT tensor-curl family \cite{Degeratu2009}.
Donnelly's equivariant calculation for the APS boundary signature operator,
as restated by Szczepański as Propositions~1--2 (explicitly attributed there
to Donnelly's Propositions~4.6--4.7)
\cite[Sec.~2, Props.~1--2]{Szczepanski2012}, gives a cotangent fixed-point
formula for a special class of affine torus actions \cite{Donnelly1978}.
This is a method antecedent for the
equivariant trigonometric route.  Those propositions neither evaluate nor
exclude the projected TT tensor-curl family considered here.
Capoferri--Vassiliev give a recent microlocal construction of curl spectral
asymmetry through regularized pseudodifferential projections
\cite{CapoferriVassiliev2025}.  This is a close analytic antecedent at
$J=1$; it does not contain the arbitrary-rank lens-space evaluation here.
The companion boundary analysis, Dirac-to-tensor-curl transfer, and bivariate
McKay--Chebyshev algebra are developed in
\cite{Mahmood2026SteinWeissAPSBoundary,
Mahmood2026DiracToTensorCurlEtaTransfer,Mahmood2026BivariateMcKay}.

Thus no novelty is attached to the operator, its sphere spectrum, cyclic
projection, or the general equivariant method.  The explicit signed evaluation
\eqref{eq:master-formula}, its composite singular resolution, and the
corollaries of Sections~6 and 7 are, to the best of our knowledge, not present
in the antecedent literature.  This is a bounded historical statement rather
than a claim of an exhaustive search.

The present result is a spectral theorem for all nontrivial round lens spaces.
It does not assemble the auxiliary eta and kernel terms of a full APS
correction, establish an off-round reducing split, or determine a complete
one-loop phase with its longitudinal, trace, ghost, framing, and counterterm
data.  Nor does the formula select a lens level or classify every eta-zero
triple $(p,J,q)$.

\section{Discussion}

The formula exhibits a useful separation of structure.  The Dedekind term is
the classical rank-one spectral asymmetry multiplied by $J$.  The residue
term records finite helicity weights, the $\delta_p$ term records modular
congruence or Fourier-aliasing coincidences, and $L_J$ repairs the finite
low-energy boundary of the physical TT spectrum.  Removable
denominator-singular character limits are handled separately by the
polynomial continuation in \eqref{eq:formal-polynomial-kernel}.  Prime and
composite levels are thereby handled by a single expression.

The most immediate mathematical continuation is a chamber decomposition of
the congruence and low-mode corrections, followed by a classification of the
remaining zero orbits.  The reflection $\eta_{p-J}=-\eta_J$ suggests that the
rank sequence should be understood as a finite sine expansion, with the
generalized Dedekind data playing the role of Fourier coefficients.  The
generating kernel \eqref{eq:cheb-generating} similarly packages the rank
dependence into a rational function, pointing toward a coefficient-extraction
approach to the zero-orbit problem.

Metric variation is also open: the projected operator is analytically
well-defined more generally, but the exact lens formula uses round quotient
representation theory.  Other natural directions include chamber formulas
for the congruence corrections and a higher-rank analogue of the rank-one
Chern--Simons normalization in Appendix~\ref{app:jeffrey}.

\appendix

\section{An equivalent Rademacher-sum form}
\label{app:rademacher}

Define
\[
 R_{p,q}(c)=\sum_{r=0}^{p-1}
 \left(\!\left(\frac rp\right)\!\right)
 \left(\!\left(\frac{qr+c}{p}\right)\!\right)
\]
and, with $S_N=0$ for $N<0$,
\[
 S_N=\sum_{a=0}^N\sum_{w\in W_a}\sum_{t\in W_{N-a}}
 R_{p,q}(wq+t).
\]

\begin{proposition}\label{prop:rademacher-form}
The eta invariant also has the form
\begin{equation}
 \eta_J=4(2S_{J-2}-S_{J-1}-S_{J-3})-L_J.
 \label{eq:rademacher-form}
\end{equation}
\end{proposition}

\begin{proof}
Let $\zeta=e^{2\pi\mathrm{i}/p}$.  The finite Fourier transform of the sawtooth
function is
\[
 \sum_{r=0}^{p-1}\left(\!\left(\frac rp\right)\!\right)
 \zeta^{-\ell r}=\frac{\mathrm{i}}{2}\cot\frac{\pi\ell}{p},
 \qquad 1\leq\ell<p.
\]
Fourier inversion and the substitution $r\mapsto qr$ therefore give
\begin{equation}\label{eq:shifted-rademacher-fourier}
 4pR_{p,q}(c)=\sum_{\ell=1}^{p-1}
 \cot\alpha_\ell\cot\beta_\ell\,\zeta^{\ell c}.
\end{equation}
Since
$\sum_{w\in W_a}\zeta^{q\ell w}=U_a(\cos2\alpha_\ell)$ and
$\sum_{t\in W_b}\zeta^{\ell t}=U_b(\cos2\beta_\ell)$, put
\[
 C_N(x,y)=\sum_{a=0}^N U_a(x)U_{N-a}(y),\qquad C_N=0\quad(N<0).
\]
The generating function of the Chebyshev polynomials gives
\[
 C_N(x,y)=\frac{U_{N+1}(x)-U_{N+1}(y)}{2(x-y)}.
\]
Using $U_{n+1}=2xU_n-U_{n-1}$ in this quotient yields
\begin{equation}\label{eq:rademacher-chebyshev-bridge}
 2C_{J-2}-C_{J-1}-C_{J-3}
 =\frac{(1-x)U_{J-1}(x)-(1-y)U_{J-1}(y)}{x-y}
 =P_J(x,y).
\end{equation}
Summing~\eqref{eq:shifted-rademacher-fourier} over the weights in $S_N$
and applying~\eqref{eq:rademacher-chebyshev-bridge} gives
$4(2S_{J-2}-S_{J-1}-S_{J-3})=\eta_J^{\mathrm{form}}$ by
\eqref{eq:formal-polynomial-kernel}.  Finally, the physical spectrum begins
at $m=J+1$, so the sub-threshold term $L_J$ is subtracted exactly as in the
proof of Theorem~\ref{thm:eta-master}.
\end{proof}

\section{Rank-one Chern--Simons normalization}
\label{app:jeffrey}

This appendix compares \eqref{eq:rank-one} with Jeffrey's spectral-flow
normalization; it is not a premise of Theorem~\ref{thm:eta-master}.
Define
\begin{equation}
 R_m(p,q)=\frac1p\sum_{k=1}^{p-1}
 \cos\left(\frac{2\pi mk}{p}\right)
 \cot\left(\frac{\pi k}{p}\right)
 \cot\left(\frac{\pi qk}{p}\right).
 \label{eq:Rm-jeffrey}
\end{equation}
Jeffrey's cotangent formula gives
\[
 R_0(p,q)=4s(q,p)=-\eta_1(L(p,q)).
\]

\begin{proposition}[Jeffrey spectral-flow control]
\label{prop:j1-cs-control}
Assume $p\geq3$ is odd and $1\leq n\leq(p-1)/2$.  Let $I_n$ be the spectral
flow between the trivial connection and the $n$th nontrivial flat $SU(2)$
connection in Jeffrey's convention.  Then
\begin{equation}
 I_n=\frac{8q^*n^2}{p}-2+R_0(p,q)-R_{2n}(p,q)
 =\frac{8q^*n^2}{p}-2-\eta_1(L(p,q))-R_{2n}(p,q).
 \label{eq:jeffrey-flow}
\end{equation}
\end{proposition}

\begin{proof}
Proposition 5.2 of \cite{Jeffrey1992} states
\[
 I_n+2=\frac{8q^*n^2}{p}
 +\frac2p\sum_{k=1}^{p-1}
 \cot\left(\frac{\pi k}{p}\right)
 \cot\left(\frac{\pi qk}{p}\right)
 \sin^2\left(\frac{2\pi nk}{p}\right).
\]
Use $2\sin^2x=1-\cos2x$ and \eqref{eq:Rm-jeffrey}.  The inverse $q^*$ is
essential.
\end{proof}

The finite Fourier identity
\[
 R_m(p,q)=4\sum_{r=0}^{p-1}
 \left(\!\left(\frac rp\right)\!\right)
 \left(\!\left(\frac{qr+m}{p}\right)\!\right),
\]
gives an exact rational evaluation of $R_m$.  Together with
$R_0=4s$ and $\eta_1=-4s$, it fixes the sign and normalization in
\eqref{eq:jeffrey-flow}.

At Chern--Simons level $\kappa$, Jeffrey's canonical two-framing invariant
also contains
\[
 \exp\left(\frac{2\pi i}{4\kappa}\,12s(q,p)\right)
 =\exp\left(\frac{2\pi i}{4\kappa}\,(-3\eta_1(L(p,q)))\right).
\]
Her comparison of the stationary-phase expression with the exact invariant is
asymptotic rather than exact at finite $\kappa$.  Accordingly, this comparison
fixes the rank-one convention but does not supply a $J\geq2$
Chern--Simons phase.

\section*{Acknowledgment of assistance}
AI-assisted research, computation, checking, and editorial support used
OpenAI Codex and Anthropic Claude. Luqman Mahmood directed the work
and is responsible for the claims and final text. These internal
checks are not external peer review.

\bibliographystyle{alpha}
\bibliography{references}

\end{document}
